\Article{25}
Before proceeding further, it will be helpful to clarify the geometrical meaning of the expressions we have obtained. To make the various directions easier to visualize, we shall follow the method we introduced in Disquiss. gen. circa superficies curvas, representing them by points on a spherical surface of radius $=1$ described about an arbitrary center. First we shall denote the directions of the coordinate axes $x,y,z$ by the points $(1)$, $(2)$, $(3)$. We then denote the direction of the normal to the surface, drawn outward with respect to the region $s$, by the point $(4)$, and finally the direction of the straight line drawn from any point of the surface toward its position after variation by the point $(5)$. For brevity we shall denote the change of position itself, or the quantity $\BareRoot{(\delta x^2+\delta y^2+\delta z^2)}$, always taken as positive, by $\delta e$. We shall write the arc between two points of the sphere, say $(1)$ and $(5)$, or the angle that measures that arc, as $(1,5)$. Thus
\[
\delta x=\delta e.\cos(1,5),\quad \delta y=\delta e.\cos(2,5),\quad \delta z=\delta e.\cos(3,5)
\]

This holds for every point of the surface. On its boundary, or perimeter $P$, two other directions come into play. First there is the direction of the element $\mathrm{d}P$, to which we assign the
\SourceLeaf{68}point $(6)$. Then there is the direction of the line perpendicular to that element, tangent to the surface and drawn inward with respect to it, to which we assign the point $(7)$. By our hypothesis the points $(6)$, $(7)$, $(4)$ lie in the same order as $(1)$, $(2)$, $(3)$. Note also that $(4,6)$, $(4,7)$, $(6,7)$ represent quadrants or right angles. This gives the equations already stated above in Article 20.
\[
\begin{gathered}
\eta Z-\zeta Y=\cos(1,7),\qquad \zeta X-\xi Z=\cos(2,7),\\
\xi Y-\eta X=\cos(3,7)
\end{gathered}
\]
The formulas of the preceding article take the following forms.
\[
\begin{aligned}
Q&=-\delta e.\cos(5,7)\\
V&=\delta e.\cos(4,5).(\frac{\mathrm{d}\xi}{\mathrm{d}x}+\frac{\mathrm{d}\eta}{\mathrm{d}y})
\end{aligned}
\]
Thus $Q$ expresses the displacement of any point of the boundary $P$ away from the plane tangent to that boundary and normal to the surface $U$, taken as positive on the side facing away from that surface. The factor of $V$ given by $\delta e.\cos(4,5)$ clearly gives the displacement of any point of the surface $U$ from its tangent plane, taken as positive on the side facing away from the region $s$.

We can also explain the other factor of $V$ by its geometrical meaning. We have
\[
\begin{gathered}
\xi=-\zeta.\frac{\mathrm{d}z}{\mathrm{d}x},\qquad \eta=-\zeta.\frac{\mathrm{d}z}{\mathrm{d}y}\\
\frac{1}{\zeta\zeta}=1+(\frac{\mathrm{d}z}{\mathrm{d}x})^2+(\frac{\mathrm{d}z}{\mathrm{d}y})^2
\end{gathered}
\]
This yields
\[
\begin{aligned}
\mathrm{d}\zeta&=\xi\zeta\zeta\mathrm{d}\frac{\mathrm{d}z}{\mathrm{d}x}+\eta\zeta\zeta\mathrm{d}\frac{\mathrm{d}z}{\mathrm{d}y}\\
\frac{\mathrm{d}\xi}{\mathrm{d}x}&=-\zeta\frac{\mathrm{dd}z}{\mathrm{d}x^2}-\frac{\mathrm{d}z}{\mathrm{d}x}.\frac{\mathrm{d}\zeta}{\mathrm{d}x}\\
&=-\zeta\frac{\mathrm{dd}z}{\mathrm{d}x^2}+\xi\xi\zeta\frac{\mathrm{dd}z}{\mathrm{d}x^2}+\xi\eta\zeta\frac{\mathrm{dd}z}{\mathrm{d}x.\mathrm{d}y}\\
&=-\zeta(\eta\eta+\zeta\zeta)\frac{\mathrm{dd}z}{\mathrm{d}x^2}+\xi\eta\zeta\frac{\mathrm{dd}z}{\mathrm{d}x.\mathrm{d}y}\\
\frac{\mathrm{d}\eta}{\mathrm{d}y}&=-\zeta\frac{\mathrm{dd}z}{\mathrm{d}y^2}+\eta\eta\zeta\frac{\mathrm{dd}z}{\mathrm{d}y^2}+\xi\eta\zeta\frac{\mathrm{dd}z}{\mathrm{d}x.\mathrm{d}y}\\
&=-\zeta(\xi\xi+\zeta\zeta)\frac{\mathrm{dd}z}{\mathrm{d}y^2}+\xi\eta\zeta\frac{\mathrm{dd}z}{\mathrm{d}x.\mathrm{d}y}
\end{aligned}
\]
and hence
\[
\begin{aligned}
\frac{\mathrm{d}\xi}{\mathrm{d}x}+\frac{\mathrm{d}\eta}{\mathrm{d}y}=-\zeta^3\big\{&\frac{\mathrm{dd}z}{\mathrm{d}x^2}[1+(\frac{\mathrm{d}z}{\mathrm{d}y})^2]-\frac{2\mathrm{dd}z}{\mathrm{d}x.\mathrm{d}y}.\frac{\mathrm{d}z}{\mathrm{d}x}.\frac{\mathrm{d}z}{\mathrm{d}y}\\
&+\frac{\mathrm{dd}z}{\mathrm{d}y^2}[1+(\frac{\mathrm{d}z}{\mathrm{d}x})^2]\big\}
\end{aligned}
\]
\SourceLeaf{69}The value of this expression is known to be
\[
=\frac{1}{R}+\frac{1}{R'}
\]
where $R,R'$ denote the extreme radii of curvature at the point in question, taken as positive whenever the surface turns its convex side outward.

\Article{26}
A careful examination of our analysis from Article 22 onward will reveal a tacit assumption underlying it. To any values of the coordinates $x,y$ there corresponds only one value of $z$, and the value of $\zeta$ is positive everywhere on the whole surface $U$. Nevertheless, the final theorem to which the analysis has led, namely (I),
\[
\delta U=-\int\delta e.\cos(5,7).\mathrm{d}P+\int\delta e.\cos(4,5).(\frac{1}{R}+\frac{1}{R'})\mathrm{d}U
\]
is valid in general and is independent of this assumption. Had we wanted to embrace this generality from the start, we would have had to encounter some complications or follow a somewhat different method. But we can also readily reach the same goal by the following considerations.

Our analysis is clearly independent of the assumption that the coordinate axis $z$ is vertical. Indeed, in it the position of the axes remains entirely arbitrary. The theorem is established for all surfaces for which the set of all points $(4)$ can be contained in a single hemisphere. It suffices to take the center (pole) of such a hemisphere as $(3)$.

If, however, we are given a surface that does not satisfy this condition, it can surely be divided into two or more parts, each of which satisfies such a condition. It will readily be seen that if a surface is divided into two parts, the validity of the theorem for the entire figure follows at once from its validity for each part. Let the figure $U$ consist of the parts $U',U''$, and let $P'$ be the boundary of the figure $U'$ and $P''$ the boundary of the figure $U''$. Further, let $P',P''$ have a common part $P'''$, so that $P'$ consists of $P'''$ and $P''''$, while $P''$ consists of $P'''$ and $P'''''$. Clearly the boundary of the whole figure $U$, namely $P$, will then consist of $P''''$ and $P'''''$. We shall indeed have
\[
\begin{aligned}
\int\delta e.\cos(5,7)\mathrm{d}P'&=\int\delta e.\cos(5,7)\mathrm{d}P'''+\int\delta e.\cos(5,7)\mathrm{d}P''''\\
\int\delta e.\cos(5,7)\mathrm{d}P''&=\int\delta e.\cos(5,7)\mathrm{d}P'''+\int\delta e.\cos(5,7)\mathrm{d}P'''''
\end{aligned}
\]
\SourceLeaf{70}but it is important to note that the value of the integral $\int\delta e.\cos(5,7)\mathrm{d}P'''$, insofar as it is part of the first integral, is exactly the opposite of the value of that same integral insofar as it is part of the second. At each point of the line $P'''$, traced in opposite directions in the two cases, the positions of the point $(7)$ are opposite and hence so are the values of the factor $\cos(5,7)$. These parts therefore cancel when we add, and we obtain
\[
\int\delta e.\cos(5,7)\mathrm{d}P'+\int\delta e.\cos(5,7)\mathrm{d}P''=\int\delta e.\cos(5,7)\mathrm{d}P
\]
Since $\delta U=\delta U'+\delta U''$, the value of $\delta U$ in agreement with the stated formula (I) follows at once when that formula is assumed to agree with the values of the variations $\delta U',\delta U''$.

Finally, we observe that the truth of theorem (I) could also have been derived from geometrical considerations, and more easily than by the analytical method. We have nevertheless preferred the latter here in order to seize the opportunity to shed some light on the calculus of variations for double integrals with variable limits, which has so far been little developed. We leave the other, geometrical method, which is fairly straightforward, to the experienced reader.

\Article{27}
It remains to work out the variations that the other elements of the expression $W$ undergo when the shape of the region $s$ is varied. We shall first deal with the variation of the volume of the region $s$.

Let us return to the two triangles considered in Article 21 and join corresponding points of their sides by straight lines to form a solid. In place of this solid we may take a prism with base $\mathrm{d}U$ and height $\xi\delta x+\eta\delta y+\zeta\delta z=\delta e.\cos(4,5)$. This expression gives a positive or negative height according as the displaced triangle, and hence the whole solid, lies outside or inside the region $s$. Hence we have (II)
\[
\delta s=\int\mathrm{d}U.\delta e.\cos(4,5)
\]
It further follows that the variation of the integral $\int z\mathrm{d}s$ is (III)
\[
\delta\int z\mathrm{d}s=\int z\mathrm{d}U.\delta e.\cos(4,5)
\]

As for the variation of the quantity $T$, we first observe that since $P$ denotes the common boundary of the surfaces $T,U$, the displacements of the
\SourceLeaf{71}points of the boundary $P$ must satisfy the condition that the new positions remain on the surface of the region $S$. Clearly, then, the displacement of the element $\mathrm{d}P$ causes the surface $T$ to undergo the change $\pm\mathrm{d}P.\delta e.\sin(5,6)$, and it is easy to see that, generally speaking, the positive or negative sign depends on the sign of the quantity $\cos(4,5)$. But this variation can be expressed more neatly by introducing a new direction in the plane tangent to the surface of the region $S$, perpendicular to the line $P$ and drawn outward with respect to the region $s$. If we denote the point corresponding to this direction by $(8)$, the variation of the surface $T$ arising from the displacement of the element $\mathrm{d}P$ will be $\mathrm{d}P.\delta e.\cos(5,8)$, or (IV)
\[
\delta T=\int\mathrm{d}P.\delta e.\cos(5,8)
\]
where the sign of the factor $\cos(5,8)$ itself determines whether the change is an increase or a decrease.

Since the point $(6)$ is the pole of the great circle drawn through the points $(7)$, $(8)$, and the point $(5)$ lies on the great circle drawn through the points $(6)$, $(8)$, the points $(5)$, $(7)$, $(8)$ form a triangle right-angled at $(8)$. Hence $\cos(5,7)=\cos(5,8).\cos(7,8)$. The arc $(7,8)$, however, measures the angle between the two planes tangent to the surfaces of the regions $s,S$ at their intersection $P$, specifically between those parts of the planes that enclose the empty space. We shall denote this angle by $i$. Then $180^0-i$ is the angle between the parts of the planes that contain the region $s$, and we have our formula (V)
\[
\cos(5,7)=\cos(5,8).\cos i.
\]

